\documentclass[english,hebrew]{article}
\usepackage[HE8,T1]{fontenc}
\usepackage[utf8]{inputenc}
\setcounter{secnumdepth}{-2}
\setlength{\parindent}{0bp}
\usepackage{amsmath}
\usepackage{amsthm}

\makeatletter

\newcommand*\LyXZeroWidthSpace{\hspace{0pt}}

\@ifundefined{date}{}{\date{}}
\usepackage{graphicx}
\usepackage{geometry}
\newcommand{\limi}{\underset{n\rightarrow\infty}{\lim}}
\newcommand{\limxz}{\underset{x\rightarrow x_{0}}{\lim}}
\newcommand{\limfxz}{\underset{x\rightarrow x_{0}}{\lim}f(x)}
\newcommand{\limfxm}{\underset{x\rightarrow x_{0}^{-}}{\lim}f(x)}
\newcommand{\limfxp}{\underset{x\rightarrow x_{0}^{+}}{\lim}f(x)}
\newcommand{\limxm}{\underset{x\rightarrow x_{0}^{-}}{\lim}}
\newcommand{\limxp}{\underset{x\rightarrow x_{0}^{+}}{\lim}}
\newcommand{\sqrm}{M_{m\times m}(\mathbb{F})}
\newcommand{\seqa}{(a_{n})_{n=1}^{\infty}}
\newcommand{\seqx}{(x_{n})_{n=1}^{\infty}}
\newcommand{\funcdr}{f:D\rightarrow\mathbb{R}}
\newcommand{\der}{\limxz\frac{f(x)-f(x_{0})}{x-x_{0}}}
\usepackage[all]{pdfprivacy}

\makeatother

\usepackage{babel}
\begin{document}
\newgeometry{top=1cm, bottom=2cm, left=1cm, right=1cm}
\title{חשבון אינפיניטסימלי {\beginL 2\endL} - אוסף נוסחאות}
\author{נעם לב}

\maketitle
\tableofcontents{}

\newpage{}

\part{שיטות אינטגרציה}

\section{נוסחאות גזירה ואינטגרציה בסיסיות}

כלל המכפלה וכלל המנה:

\L{%
\[
(u\cdot v)'=u'v+uv',\quad\left(\frac{u}{v}\right)'=\frac{u'v-uv'}{v^{2}}
\]
}%
כלל השרשרת:\L{%
\[
\left(f(g(x))\right)'=f'(g(x))\cdot g'(x)
\]
}%

סינוס וקוסינוס:

\L{%
\[
(\sin(x))\prime=\cos(x),\quad(\cos(x))\prime=-\sin(x)
\]
}%

טנגנס, קוטנגנס, ארקטנגנס:

\L{%
\[
(\tan(x))\prime=\frac{1}{\cos^{2}(x)}\LyXZeroWidthSpace=1+\tan^{2}(x),\quad(\cot x)\prime=-\frac{1}{\sin^{2}(x)},\quad(\arctan x)'=\frac{1}{1+x^{2}}\LyXZeroWidthSpace
\]
}%

ארקסינוס וארקקוסינוס:\L{%
\[
(\arcsin(x))'=\frac{1}{\sqrt{1-x^{2}}},\quad(\arccos(x))'=-\frac{1}{\sqrt{1-x^{2}}}
\]
}%

אינטגרל של פונקציית חזקה:\L{%
\[
\int x^{n}\,dx=\frac{x^{n+1}}{n+1}+C
\]
}%

אינטגרל של אחד חלקי המשתנה ונגזרת פנימית במונה:\L{%
\[
\int\frac{1}{x}\,dx=\ln|x|+C,\quad\int\frac{f'(x)}{f(x)}\,dx=\ln|f(x)|+C
\]
}%

\section{אינטגרציה בחלקים}

אפשר להיזכר בה בעזרת כלל המכפלה.\L{%
\[
(u\cdot v)'=u'v+uv'
\]
}%

לכן\L{%
\[
\int(u(x)\cdot v(x))'\ dx=\int\left(u'(x)v(x)+u(x)v'(x)\right)\ dx\Rightarrow u(x)\cdot v(x)=\int\left(u'(x)v(x)+u(x)v'(x)\right)\ dx
\]
}%
\L{%
\[
\underset{\text{linearity}}{\Rightarrow}u(x)\cdot v(x)=\int u(x)'v(x)\ dx+\int u(x)v'(x)\ dx
\]
}%

לכן לאחר העברת אגפים\L{%
\[
\int u(x)v'(x)\,dx=u(x)v(x)-\int u'(x)v(x)\,dx
\]
}%

עדיף לקרוא \L{$v'$} למישהו שנוח לעשות לו אינטגרל, ו-\L{$u$} למשהו
שנוח לגזור. 

דהיינו, בשביל להשתמש בשיטה זו, נזהה באינטגרל מכפלה של ביטוי שקל לגזור
)למשל, אם גזירה מורידה את המעלה שלו( וביטוי שקל לטנגרל בלי לסבך את
המבנה שלו.

אם יש פונקציה יחידה, נתייחס אליה כמכפלה בקבוע )אפילו \L{$1$}(.

\section{אינטגרציה בהצבה}

\L{%
\[
\int f(\varphi(x))\cdot\varphi'(x)\,dx=F(\varphi(x))+C
\]
}%

אפשר להיזכר בה בעזרת כלל השרשרת.

\L{%
\[
\left(f(g(x))\right)'=f'(g(x))\cdot g'(x)
\]
}%
לכן\L{%
\[
f(g(x))+C=\int f'(g(x))\cdot g'(x)\ dx
\]
}%

)לכתיבה הריגורוזית ראו את הנוסחה הראשונה(.

שיטה זו נוחה לשילוב עם כתיב לייבניץ לנגזרת )בוחרים \L{$u=\varphi(x)$}
ומטנגרלים לפי \L{$u$}(. 

\section{כתיב לייבניץ לנגזרת}

המשמעות של \L{$\int f(x)\ dx$} היא שמטנגרלים את הפונקציה \L{$f(x)$}
לפי המשתנה \L{$x$}.

לעיתים, נרצה לבצע החלפת משתנה, ולטנגרל אותה לפי משתנה אחר.

נרצה לבחור משתנה \L{$u$} שמצד אחד עוזר לנו לפתור את השאלה )למשל ביטוי
שהנגזרת שלו כבר קיימת באינטגרל( אך מצד שני לא מסבך את המבנה כשגוזרים
אותו.

לאחר מכן, נחשב את \L{$dx$} בעזרת \L{$du$}: \L{$\frac{du}{u'}=dx$}. 

נחליף באינטגרל את \L{$dx$} ו-\L{$x$} בהתאם, נפתור, ונציב בחזרה את
\L{$u$}.

\newpage{}

\part{זהויות טריגונומטריות}

זהויות זווית כפולה הן זהויות שחשוב לזכור:\L{%
\[
1.\quad\cos\left(2x\right)=2\cos^{2}(x)-1
\]
}%
\L{%
\[
2.\quad\sin(2x)=2\sin(x)\cos(x)
\]
}%

לפי משפט פיתגורס, אפשר לקבל ש-\L{%
\[
3.\quad\sin^{2}(x)+\cos^{2}(x)=1
\]
}%

בעזרת זהויות אלו, אפשר לקבל זהות שימושית של הורדת דרגה:\L{%
\[
4.\quad\cos\left(2x\right)=2\cos^{2}(x)-1\underset{(1)}{\Rightarrow}\cos^{2}(x)=\frac{\cos(2x)+1}{2}
\]
}%
\L{%
\[
5.\quad\Rightarrow\sin^{2}(x)\underset{4\text{ and }3}{=}\frac{1-\cos(2x)}{2}
\]
}%

בנוסף, אפשר לחלק את שני האגפים של זהות שלוש ב-\L{$\cos^{2}(x)$} ובעזרת
ההגדרה \L{$\tan(x)=\frac{\sin(x)}{\cos(x)}$} לקבל

\L{%
\[
6.\quad\tan^{2}(x)+1=\frac{1}{\cos^{2}(x)}
\]
}%

סכומים בסינוס וקוסינוס:\L{%
\[
7.\quad\sin(\alpha\pm\beta)=\sin(\alpha)\cos(\beta)\pm\cos(\alpha)\sin(\beta)
\]
}%
\L{%
\[
8.\quad\cos(\alpha\pm\beta)=\cos(\alpha)\cos(\beta)\mp\sin(\alpha)\sin(\beta)
\]
}%

\end{document}
